\section{Disponibilidad}

\begin{caso}{Intrusión y eliminación no autorizada de datos bancarios (Corea del Sur, Shinhan Bank)}
Leer el caso de intrusión y eliminación no autorizada de datos bancarios.
\end{caso}

\begin{fuentes}
  \fuente{https://www.seeker.com/massive-data-deleting-attack-hits-south-korea-1767326593.html}
  \fuente{https://www.nbcnews.com/id/wbna51256632}
  \fuente{https://www.darkreading.com/cyberattacks-data-breaches/south-korean-banks-lose-data-in-malware-attack}
  \item \footnotesize Otros disponibles en la web.
\end{fuentes}
\nocite{sk-seeker,sk-nbc,sk-darkreading}

\pregunta{¿Cuál era el objetivo final del atacante?}
\begin{respuesta}
\end{respuesta}

\pregunta{¿Cómo se logró la eliminación de los datos o servicios?}
\begin{respuesta}
\end{respuesta}

\pregunta{¿Qué mecanismo(s) pudo haberse implementado para evitar la eliminación no
autorizada de datos?}
\begin{respuesta}
\end{respuesta}

\subsection*{Validación de entendimiento del caso}

\begin{mcq}{En el caso de Shinhan Bank estudiado previamente, ¿cuál fue el vector de
entrada más probable utilizado por el atacante? (1 o más respuestas válidas):}
  \opcion{Explotación de una vulnerabilidad en el portal web de Shinhan Bank que
    permitió afectar el sistema operativo de las máquinas corporativas.}
  \opcion{Equipos corporativos que hacían parte de una botnet y a través de los cuales
    fue posible la eliminación de archivos de sistema operativo.}
  \opcion{Inserción de una USB infectada por parte de uno de los empleados, a través de
    la cual se realizó la implantación de un malware de corrupción del sistema operativo.}
\end{mcq}
